\section{Conclusiones}

Tras el desarrollo experimental y analítico de esta práctica, se logró validar la utilidad de los teoremas de redes como herramientas de simplificación en el análisis de circuitos complejos. La correlación entre los datos obtenidos y las simulaciones permitió confirmar que, más allá de la teoría, estos principios rigen el comportamiento físico de los sistemas lineales y deterministas.


En primer lugar, el \textbf{Teorema de Sustitución} demostró ser una técnica poderosa para aislar secciones de una red. Al reemplazar una rama por una fuente de voltaje equivalente, comprobamos que las condiciones del resto del circuito permanecen invariantes. El hecho de que las mediciones de $V_{R1}$ a $V_{R6}$ se mantuvieran constantes tras la modificación valida el concepto de unicidad de la solución, de esta manera sabemos que si el voltaje en los nodos se preserva, la respuesta del sistema es única.


Respecto al \textbf{Teorema de Tellegen}, los resultados fueron contundentes en cuanto a la conservación de la energía. Aunque en la práctica obtuvimos una diferencia de apenas $0.3$ mW entre la potencia suministrada y la absorbida, este margen de error es totalmente aceptable en un entorno experimental. Esto nos recuerda que, a pesar de las tolerancias de las resistencias y la precisión de los multímetros, la suma algebraica de las potencias en una red siempre tiende a cero.


Finalmente, el \textbf{Teorema de Thevenin} resultó ser el más revelador en cuanto a la práctica de ingeniería. La mayor dificultad radicó en la calibración de las fuentes para el circuito equivalente; sin embargo, incluso con esa ligera desviación respecto a los $-4.26$ V medidos inicialmente, el modelo predijo correctamente el comportamiento de la red ante cambios de carga.


De esta práctica se puede concluir que dominar estos teoremas no solo facilita los cálculos teóricos, sino que permite entender la infraestructura de redes más grandes mediante modelos simplificados.